\documentclass[11pt,fontset=fandol]{ctexart}
\usepackage[a4paper,margin=2.2cm]{geometry}
\usepackage{amsmath,amssymb,enumitem,xcolor,tcolorbox}
\tcbuselibrary{breakable}
\setlength{\parindent}{0pt}
\setlength{\parskip}{0.42em}
\newcommand{\R}{\mathbb R}
\newcommand{\answer}{\par\textbf{解答}\quad}
\newcommand{\basis}{\par\textbf{每一步凭什么}\quad}
\newcommand{\reaction}{\par\textbf{必须形成的反应}\quad}
\newenvironment{drill}[2]{
  \begin{tcolorbox}[colback=white,colframe=black!35,title={#1\quad\small #2},breakable]
}{\end{tcolorbox}}

\title{导数基础短逻辑题}
\author{}
\date{}

\begin{document}
\maketitle

这套题从甲种本《微积分初步》、人教B版选择性必修三和文件夹中的导数题库中选题。题目不是为了凑数量，也不把完整导数大题切成纯计算题；每题只留下一个在大题中会反复出现、必须低延迟调用的逻辑动作。

\begin{drill}{1. 极限里的增量不是原封不动的}{甲种本《微积分初步》导数定义题型改编}
已知 $f$ 在 $x_0$ 处可导，且
\[
\lim_{\Delta x\to0}\frac{f(x_0+3\Delta x)-f(x_0)}{\Delta x}=6.
\]
求 $f'(x_0)$。

\answer 分母中的 $3\Delta x$ 才是自变量的增量。令 $h=3\Delta x$，则
\[
6=\lim_{\Delta x\to0}3\cdot\frac{f(x_0+3\Delta x)-f(x_0)}{3\Delta x}
=3f'(x_0).
\]
所以 $f'(x_0)=2$。

\basis 导数定义中的分母必须是函数输入真正增加的量。这里输入从 $x_0$ 变到 $x_0+3\Delta x$，增加的是 $3\Delta x$。

\reaction 看见极限长得像导数定义但分母和括号内的增量不一致，先找“真正增加了多少”，再把多出来的常数提出来。
\end{drill}

\begin{drill}{2. 反向写的差商会改变符号}{甲种本《微积分初步》导数定义题型改编}
已知
\[
\lim_{\Delta x\to0}\frac{f(2)-f(2-\Delta x)}{\Delta x}=-3.
\]
求 $f'(2)$。

\answer 令 $h=-\Delta x$，则 $h\to0$，且
\[
\frac{f(2)-f(2-\Delta x)}{\Delta x}
=\frac{f(2)-f(2+h)}{-h}
=\frac{f(2+h)-f(2)}{h}.
\]
所以 $f'(2)=-3$。

\basis 分子和分母同时反向，两个负号抵消。不能只看到 $2-\Delta x$ 就凭感觉多添一个负号。

\reaction 差商不是标准样子时，把“新输入减旧输入”写清楚；分子、分母的方向要一起看。
\end{drill}

\begin{drill}{3. 导数就是切线斜率}{甲种本《微积分初步》切线习题改编}
曲线 $y=f(x)$ 在点 $(1,2)$ 处的切线经过点 $(3,8)$。求 $f'(1)$。

\answer 切线斜率为
\[
\frac{8-2}{3-1}=3.
\]
而该斜率就是 $f'(1)$，故 $f'(1)=3$。

\basis 导数 $f'(x_0)$ 的几何意义是曲线在 $(x_0,f(x_0))$ 处切线的斜率。

\reaction 题目给出“切线经过两个点”，先别急着求导；两点坐标已经直接给了导数值。
\end{drill}

\begin{drill}{4. 切线方程需要两个信息}{甲种本《微积分初步》切线习题改编}
已知 $f(2)=5,\ f'(2)=-4$，求曲线 $y=f(x)$ 在 $x=2$ 处的切线方程。

\answer 切点是 $(2,5)$，斜率是 $-4$，所以
\[
y-5=-4(x-2).
\]

\basis 一条直线由一个点和斜率确定；这里点来自 $f(2)$，斜率来自 $f'(2)$。

\reaction 求切线方程时强制写两行：切点 $(x_0,f(x_0))$；斜率 $f'(x_0)$。缺一项就不能写点斜式。
\end{drill}

\begin{drill}{5. 与 $x$ 轴平行翻译成导数为零}{项目导数题库改编}
曲线 $y=x^3-3x+1$ 在点 $P$ 处的切线与 $x$ 轴平行。求点 $P$ 的横坐标。

\answer 切线与 $x$ 轴平行，斜率为 $0$。而
\[
y'=3x^2-3.
\]
故
\[
3x^2-3=0,
\]
所以 $x=\pm1$。

\basis 水平直线的斜率为 $0$；曲线在某点的切线斜率就是该点的导数值。

\reaction “切线平行于 $x$ 轴”立刻翻译成 $f'(x_0)=0$；求出横坐标后不要忘记题目问的是横坐标还是完整坐标。
\end{drill}

\begin{drill}{6. 垂直不是斜率相等}{项目导数题库改编}
曲线 $y=x^2+\ln x$ 在点 $x=x_0$ 处的切线垂直于直线 $2x+y-1=0$。求 $x_0$。

\answer 已知直线斜率为 $-2$，所以切线斜率为 $\frac12$。又
\[
f'(x)=2x+\frac1x,
\]
故
\[
2x_0+\frac1{x_0}=\frac12.
\]
这在 $x_0>0$ 上无解，因为 $2x_0+\frac1{x_0}\ge2\sqrt2>\frac12$。所以不存在这样的点。

\basis 两条非竖直直线垂直时斜率乘积为 $-1$，不是斜率相等。

\reaction 看见“平行”才令斜率相等；看见“垂直”先把已知直线斜率取负倒数，同时把函数定义域带上。
\end{drill}

\begin{drill}{7. 常数项求导后消失}{甲种本《微积分初步》求导习题改编}
已知曲线
\[
y=x^3+ax^2+bx+7
\]
在 $x=1$ 处的切线斜率为 $5$。写出 $a,b$ 满足的关系。

\answer
\[
y'=3x^2+2ax+b.
\]
代入 $x=1$：
\[
3+2a+b=5,
\]
即 $2a+b=2$。

\basis 切线斜率只由导函数在该点的值决定；常数 $7$ 的导数为 $0$，不会进入斜率条件。

\reaction 题目只给“斜率”时，先求导再代横坐标；不要把切点函数值条件擅自加进去。
\end{drill}

\begin{drill}{8. 乘积不能分别求导后相乘}{甲种本《微积分初步》乘法法则习题改编}
求函数
\[
f(x)=x^2\ln x\quad(x>0)
\]
的导数。

\answer 把它看作两个函数的乘积：
\[
f'(x)=(x^2)'\ln x+x^2(\ln x)'
=2x\ln x+x.
\]

\basis $(uv)'=u'v+uv'$；只有乘积中有一个因子是常数时，才会退化成“常数乘导数”。

\reaction 看见两个会随 $x$ 变化的部分相乘，先在脑中标出 $u,v$，然后强制写 $u'v+uv'$。
\end{drill}

\begin{drill}{9. 商法则的分母要整体平方}{甲种本《微积分初步》除法法则习题改编}
求函数
\[
f(x)=\frac{x+1}{x-1}
\]
的导数。

\answer
\[
f'(x)=\frac{1\cdot(x-1)-(x+1)\cdot1}{(x-1)^2}
=-\frac2{(x-1)^2}.
\]

\basis $(u/v)'=(u'v-uv')/v^2$。分母平方的是整个分母 $v$，不是只平方其中某一项。

\reaction 分式求导时按固定顺序写：分子导数乘分母，减分子乘分母导数，最后除以分母平方。别靠记忆口算符号。
\end{drill}

\begin{drill}{10. 复合函数先认内层}{甲种本《微积分初步》复合函数习题改编}
求函数
\[
f(x)=(3x-1)^5
\]
的导数。

\answer 把 $3x-1$ 当成整体：
\[
f'(x)=5(3x-1)^4\cdot3=15(3x-1)^4.
\]

\basis 外层 $u^5$ 对 $u$ 求导得到 $5u^4$；但 $u=3x-1$ 仍随 $x$ 变化，还要乘内层导数 $u'=3$。

\reaction 幂的底数不是单独的 $x$ 时，先问“底数这个整体对 $x$ 的导数是多少”，这就是链式法则。
\end{drill}

\begin{drill}{11. 指数的底数和指数不要混淆}{人教B版初等函数导数题型改编}
求函数
\[
f(x)=2^{x^2}
\]
的导数。

\answer
\[
f'(x)=2^{x^2}\ln2\cdot2x=2x\ln2\cdot2^{x^2}.
\]

\basis $(a^u)'=a^u\ln a\cdot u'$；这里底数是常数 $2$，指数是内层 $x^2$。

\reaction 看见 $a^{u(x)}$，把 $\ln a$ 和 $u'(x)$ 一起带出来；不要把它错当成 $(x^2)^2$ 一类幂函数。
\end{drill}

\begin{drill}{12. 对数复合也要乘内层导数}{人教B版初等函数导数题型改编}
求函数
\[
f(x)=\ln(2x+1)
\]
的导数，并写出定义域。

\answer 定义域为 $2x+1>0$，即 $x>-\frac12$。在该定义域内，
\[
f'(x)=\frac{2}{2x+1}.
\]

\basis $(\ln u)'=u'/u$；先有 $u=2x+1>0$，对数函数才存在。

\reaction 含 $\ln(\text{式子})$ 的题同时触发两件事：先写真数大于零；求导时写“内层导数除以内层”。
\end{drill}

\begin{drill}{13. 定义求导时，先把分子化出增量}{甲种本《微积分初步》导数习题改编}
用导数定义求
\[
f(x)=\sqrt{x+4}
\]
在 $x=5$ 处的导数。

\answer
\[
f'(5)=\lim_{\Delta x\to0}\frac{\sqrt{9+\Delta x}-3}{\Delta x}.
\]
分子有根式相减，乘共轭式：
\[
\frac{\sqrt{9+\Delta x}-3}{\Delta x}
=\frac{1}{\sqrt{9+\Delta x}+3}.
\]
因此
\[
f'(5)=\frac16.
\]

\basis 导数定义中不能直接把 $\Delta x=0$ 代进一个分母含 $\Delta x$ 的式子。先把分子中的 $\Delta x$ 消掉，才能取极限；根式相减正是乘共轭式的信号。

\reaction 用定义求导时，分式先看“怎样把分子化出 $\Delta x$”。遇到根式相减，直接尝试乘共轭式。
\end{drill}

\begin{drill}{14. 与已知直线平行，先翻译成斜率相等}{甲种本《微积分初步》习题五改编}
曲线
\[
y=x^3+3x
\]
在横坐标为 $x_0$ 的点处的切线与直线 $y=15x+2$ 平行。求 $x_0$。

\answer 已知直线的斜率为 $15$。曲线在 $x_0$ 处的切线斜率为
\[
f'(x_0)=3x_0^2+3.
\]
所以
\[
3x_0^2+3=15,
\]
即
\[
x_0=\pm2.
\]

\basis 两直线平行时斜率相等；曲线切线的斜率由导数给出。题目还没要求切线方程，所以到横坐标已经完成。

\reaction 看见“曲线的切线与直线平行”，先把直线斜率取出来，令 $f'(x_0)$ 等于它；不要一开始就写切线方程。
\end{drill}

\begin{drill}{15. 切线倾斜角也只是一个斜率条件}{甲种本《微积分初步》习题五改编}
曲线
\[
y=x^3+x^2-1
\]
在哪些横坐标处的切线倾斜角为 $45^\circ$？

\answer 倾斜角为 $45^\circ$ 的直线斜率为
\[
\tan45^\circ=1.
\]
又
\[
f'(x)=3x^2+2x.
\]
因此
\[
3x^2+2x=1,
\]
解得
\[
x=-1\quad\text{或}\quad x=\frac13.
\]

\basis 曲线的切线倾斜角不是另一个新对象，它仍然通过“切线斜率 $=\tan\alpha$”翻译成导数值。

\reaction 看见切线倾斜角，先翻译成 $f'(x_0)=\tan\alpha$。这里问的是横坐标，不必再把它们代回函数求纵坐标。
\end{drill}

\begin{drill}{16. $f'(x_0)=f(x_0)$ 是一条可以直接展开的方程}{甲种本《微积分初步》习题五改编}
已知
\[
f(x)=x^2(x-1),
\]
求满足 $f'(x_0)=f(x_0)$ 的 $x_0$。

\answer
\[
f'(x)=3x^2-2x.
\]
所以
\[
3x_0^2-2x_0=x_0^2(x_0-1).
\]
移项、提取公因式：
\[
x_0(x_0^2-4x_0+2)=0.
\]
故
\[
x_0=0,\quad 2-\sqrt2,\quad 2+\sqrt2.
\]

\basis $f'(x_0)$ 和 $f(x_0)$ 都是数，题目已经给出了它们相等的条件。先把导函数算出来，再把两个表达式放进同一个方程处理即可。

\reaction 题目直接给出“导数值 $=$ 函数值”这类关系时，不要试图猜图像。先按定义把 $f'(x)$ 写出，再代 $x_0$ 列方程。
\end{drill}

\begin{drill}{17. 可导奇函数的导函数是偶函数}{人教B版导数与函数性质题型改编}
已知可导函数 $f$ 是奇函数，且 $f'(2)=3$。求 $f'(-2)$。

\answer 因为 $f$ 是可导奇函数，所以 $f'$ 是偶函数，故
\[
f'(-2)=f'(2)=3.
\]

\basis $f(-x)=-f(x)$ 两边对 $x$ 求导，左边由链式法则得到 $-f'(-x)$，于是 $-f'(-x)=-f'(x)$。

\reaction 题目同时出现“可导、奇偶性、导函数在某点的值”，先转换到导函数的奇偶性：奇函数导数偶，偶函数导数奇。
\end{drill}

\begin{drill}{18. 导数为零并不自动等于极值}{甲种本《微积分初步》极值章节例题改编}
函数 $f(x)=x^3$ 在 $x=0$ 处满足 $f'(0)=0$。问 $x=0$ 是否为极值点？

\answer 不是。因为
\[
f'(x)=3x^2>0\quad(x\ne0),
\]
在 $0$ 左右导数都为正，没有变号，所以函数在两侧都递增。

\basis 极值要求函数在该点左右的增减趋势不同；$f'(x_0)=0$ 只是可疑点，不是极值的充分条件。

\reaction 解出 $f'(x)=0$ 后不能直接写“极值点”。下一步必须看导数在该点左右的符号有没有变化。
\end{drill}

\begin{drill}{19. 导数不存在也可能是极值点}{甲种本《微积分初步》可导性章节例题改编}
函数 $f(x)=|x|$ 在 $x=0$ 处是否有极值？导数是否存在？

\answer $x=0$ 是极小值点，极小值为 $0$；但 $f'(0)$ 不存在。

\basis 当 $x<0$ 时 $f(x)=-x$，斜率为 $-1$；当 $x>0$ 时 $f(x)=x$，斜率为 $1$。左右导数不同，所以不可导；但附近函数值都大于 $f(0)$。

\reaction “极值点一定满足导数为零”是错的。先问函数在左右两边怎样变化；不可导点也应被列入极值候选。
\end{drill}

\begin{drill}{20. 导数从正到负是极大值}{甲种本《微积分初步》极值章节改编}
已知函数 $f$ 在 $x=2$ 附近可导，且
\[
f'(x)>0\ (x<2),\qquad f'(x)<0\ (x>2).
\]
判断 $x=2$ 处的极值类型。

\answer $x=2$ 是极大值点。

\basis 左边导数为正说明从左向右函数递增；右边导数为负说明继续向右函数递减。因此函数先上升后下降，顶点在 $x=2$。

\reaction 不要背“极大极小”的字面。直接读导数符号：$+\to-$ 是先升后降，极大；$-\to+$ 是先降后升，极小。
\end{drill}

\begin{drill}{21. 区间最值不能漏端点}{甲种本《微积分初步》最值章节改编}
求
\[
f(x)=x^3-3x
\]
在 $[-2,2]$ 上的最大值和最小值。

\answer
\[
f'(x)=3x^2-3=0\Longrightarrow x=\pm1.
\]
比较候选点：
\[
f(-2)=-2,\quad f(-1)=2,\quad f(1)=-2,\quad f(2)=2.
\]
所以最大值为 $2$，最小值为 $-2$。

\basis 闭区间连续函数的最值只可能在区间端点或区间内部极值点取得。

\reaction 闭区间最值题先列候选点清单：左端点、右端点、区间内部的导数为零或不可导点。少一个就可能漏答案。
\end{drill}

\begin{drill}{22. 驻点不在区间内不能参与比较}{项目导数题库改编}
求
\[
f(x)=x^3-3x
\]
在 $[1,3]$ 上的最小值。

\answer $f'(x)=3x^2-3=0$ 给出 $x=\pm1$。区间内部没有导数为零的点，$x=1$ 只是端点。比较端点：
\[
f(1)=-2,\qquad f(3)=18.
\]
所以最小值为 $-2$。

\basis 最值候选中的极值点必须落在开区间 $(1,3)$ 内；端点本来会单独比较。

\reaction 解出导数零点后，先和题目给的区间取交。不要把定义域外、区间外的点带进表格。
\end{drill}

\begin{drill}{23. 导函数符号表的分界点来自定义域和零点}{步步高导数单调性题型改编}
求函数
\[
f(x)=x+\frac4x-3\ln x
\]
的单调递减区间。

\answer 定义域为 $(0,+\infty)$。求导：
\[
f'(x)=1-\frac4{x^2}-\frac3x
=\frac{(x-4)(x+1)}{x^2}.
\]
在 $x>0$ 上，$x^2>0,\ x+1>0$，故 $f'(x)<0$ 等价于 $x<4$。所以递减区间为
\[
(0,4).
\]

\basis 导数符号只在函数定义域内讨论；在定义域内能确定恒正的因子可以不再反复分类。

\reaction 做导数符号题的第一行写定义域。分解后先划掉在定义域内恒正的部分，把真正决定符号的因子留下。
\end{drill}

\begin{drill}{24. 导数是分式时，先把分子做成明显非负}{甲种本《微积分初步》单调性习题改编}
证明函数
\[
f(x)=x+\frac1{x^2+1}
\]
在 $\R$ 上单调递增。

\answer
\[
f'(x)=1-\frac{2x}{(x^2+1)^2}
=\frac{x^4+2x^2-2x+1}{(x^2+1)^2}
=\frac{x^4+x^2+(x-1)^2}{(x^2+1)^2}>0.
\]
故 $f$ 在 $\R$ 上单调递增。

\basis 这里分母恒正，导数的正负只由分子决定；把分子拆成若干个平方项的和，就能把“恒正”写出来。

\reaction 导函数是分式时，先处理分母的符号。剩下的分子若不容易看正负，就尝试配方、拆平方，把“恒非负”写出来。
\end{drill}

\begin{drill}{25. 函数递增不能反推导数处处大于零}{步步高导数单调性辨析题}
判断命题：“函数 $f$ 在区间 $I$ 上单调递增，则对任意 $x\in I$ 都有 $f'(x)>0$。”

\answer 命题错误。反例：$f(x)=x^3$ 在 $\R$ 上严格递增，但
\[
f'(0)=0.
\]

\basis $f'(x)>0$ 是推出严格递增的充分条件之一，不是严格递增的必要条件。

\reaction 遇到“单调递增 $\Rightarrow$ 导数恒正”这种反向说法，立刻想到 $x^3$：它递增，但在 $0$ 处斜率为零。
\end{drill}

\begin{drill}{26. 用导数研究单调性前先看定义域}{步步高导数单调性题型改编}
求函数
\[
f(x)=\ln x-\frac1x
\]
的单调区间。

\answer 定义域为 $(0,+\infty)$，且
\[
f'(x)=\frac1x+\frac1{x^2}=\frac{x+1}{x^2}>0.
\]
所以 $f$ 在 $(0,+\infty)$ 上单调递增。

\basis $\ln x$ 和 $1/x$ 都要求 $x>0$；只有在这个范围内，$x^2$ 才能视为正。

\reaction 有对数、分式、根式时，定义域不是形式步骤。它决定后面符号判断时哪些量可以认定为正。
\end{drill}

\begin{drill}{27. 含参单调性先把导数变成关于 $x$ 的不等式}{项目导数题库改编}
若函数
\[
f(x)=x^2-2ax
\]
在 $[1,3]$ 上单调递增，求 $a$ 的取值范围。

\answer
\[
f'(x)=2x-2a.
\]
要在 $[1,3]$ 上递增，需要对任意 $x\in[1,3]$ 有
\[
2x-2a\ge0,
\]
即 $a\le x$ 对任意 $x\in[1,3]$ 成立。因此 $a$ 必须不超过这段 $x$ 的最小值：
\[
a\le1.
\]

\basis “对任意 $x$，$a\le x$”不是代一个随便的 $x$，而是要求 $a$ 压在所有 $x$ 的下方。

\reaction 含参单调性先求导，把题目翻译成“导函数在整个区间的正负”。参数孤立后，再按任意/存在翻译成最值。
\end{drill}

\begin{drill}{28. 导数符号与参数的端点效应}{项目导数题库改编}
若函数
\[
f(x)=x^2-2ax
\]
在 $[1,3]$ 上单调递减，求 $a$ 的取值范围。

\answer 要求
\[
f'(x)=2x-2a\le0\quad(x\in[1,3]),
\]
即 $a\ge x$ 对任意 $x\in[1,3]$ 成立。因此
\[
a\ge3.
\]

\basis 要压住整个区间内所有的 $x$，就要压住其中最大的那个。

\reaction 同一个导函数不等式，参数在左边还是右边会决定你看最小值还是最大值。别只记上一题的答案方向。
\end{drill}

\begin{drill}{29. 极值点条件先给出参数方程}{项目导数题库改编}
已知
\[
f(x)=x^3-3ax+1
\]
在 $x=1$ 处有极值，求 $a$。

\answer 极值点是导数为零的必要条件：
\[
f'(x)=3x^2-3a,\qquad f'(1)=3-3a=0.
\]
所以 $a=1$。此时 $f'(x)=3(x^2-1)$，在 $1$ 左右由负变正，确实有极小值。

\basis 若函数在内点可导且有极值，则该点导数必为零；求出参数后仍应验证导数变号，确认它不是驻点。

\reaction 题目说“$x_0$ 是极值点”，第一步永远是 $f'(x_0)=0$；这只是在求候选参数，最后还要看左右符号。
\end{drill}

\begin{drill}{30. 两个极值点意味着导函数至少有两个变号零点}{项目导数题库改编}
若
\[
f(x)=x^3+ax^2+bx
\]
有两个不同的极值点，写出关于 $a,b$ 的必要且充分条件。

\answer
\[
f'(x)=3x^2+2ax+b.
\]
二次函数要有两个不同实根，且二次项系数非零，需
\[
\Delta=(2a)^2-12b>0,
\]
即
\[
a^2>3b.
\]
两个单根处导数符号都会改变，所以此条件也充分。

\basis 三次函数导数是开口向上的二次函数；两个不同实根把数轴分成 $+,-,+$ 三段，函数就有一个极大值和一个极小值。

\reaction “有两个极值点”先翻译成“导函数有两个不同零点并变号”。对三次函数，这通常就是导函数二次方程判别式大于零。
\end{drill}

\begin{drill}{31. 已知极值点横坐标，可用导函数因式分解}{项目导数题库改编}
已知三次函数 $f$ 的导函数是首项系数为 $3$ 的二次函数，且 $x=1,x=4$ 是 $f$ 的两个极值点。写出 $f'(x)$。

\answer 导函数的两个零点为 $1,4$，首项系数为 $3$，所以
\[
f'(x)=3(x-1)(x-4).
\]

\basis 一个二次函数由首项系数和两个根唯一确定。

\reaction 已知两个极值点的横坐标时，不必把导函数写成一般式再列方程；直接写成“首项系数 $\times$ 两个一次因子”。
\end{drill}

\begin{drill}{32. 切线经过原点给出的是 $f(x_0)=x_0f'(x_0)$}{项目导数题库改编}
曲线 $y=f(x)$ 在点 $(x_0,f(x_0))$ 处的切线经过原点。写出 $f(x_0)$ 与 $f'(x_0)$ 的关系。

\answer 切线方程为
\[
y-f(x_0)=f'(x_0)(x-x_0).
\]
代入原点 $(0,0)$：
\[
-f(x_0)=-x_0f'(x_0),
\]
故
\[
f(x_0)=x_0f'(x_0).
\]

\basis “切线经过某点”就是把该点坐标代入切线方程。

\reaction 看见切线经过固定点，不要先画图猜。先写切线点斜式，再把固定点代进去；这是把几何条件变成代数关系。
\end{drill}

\begin{drill}{33. 求参数切线时，函数值和导数值是两条条件}{项目导数题库改编}
直线 $y=2x+1$ 是曲线
\[
y=x^2+ax+b
\]
在 $x=1$ 处的切线，求 $a,b$。

\answer 切点在直线上也在曲线上：
\[
1+a+b=3.
\]
切线斜率相等：
\[
2+a=2.
\]
故 $a=0,\ b=2$。

\basis “是 $x=1$ 处的切线”包含两层：曲线经过 $(1,3)$，且曲线在此处的斜率是 $2$。

\reaction 给出一条具体直线作为切线时，默认建立两条方程：函数值相等、导数值等于直线斜率。
\end{drill}

\begin{drill}{34. 切点未知时不要只列斜率条件}{项目导数题库改编}
求曲线
\[
y=x^2
\]
与直线 $y=2x+c$ 相切时 $c$ 的值。

\answer 设切点横坐标为 $x_0$。斜率相等：
\[
2x_0=2\Longrightarrow x_0=1.
\]
切点在直线上：
\[
1=2+c,
\]
故 $c=-1$。

\basis 两曲线相切不仅斜率相同，而且在同一个点有相同函数值。

\reaction “相切”不能只让导数相等。切点未知时设 $x_0$，同时写函数值相等和导数值相等。
\end{drill}

\begin{drill}{35. 证明不等式先把目标移成函数}{项目导数题库构造题改编}
证明当 $x\ge0$ 时，
\[
e^x\ge x+1.
\]

\answer 令
\[
g(x)=e^x-x-1.
\]
则
\[
g'(x)=e^x-1\ge0\quad(x\ge0),
\]
所以 $g$ 在 $[0,+\infty)$ 上递增，且
\[
g(x)\ge g(0)=0.
\]
原不等式成立。

\basis 要证明左边大于右边，就研究它们的差是否非负；导数用来证明这个差函数从一个已知点开始怎样变化。

\reaction 证明 $A(x)\ge B(x)$ 时，先主动问“$A(x)-B(x)$ 的正负能不能用单调性解决”。这是构造函数的最基本起手式。
\end{drill}

\begin{drill}{36. 构造函数后，锚点要选能算出零的点}{项目导数题库构造题改编}
证明当 $x>0$ 时，
\[
\ln x\le x-1.
\]

\answer 令
\[
g(x)=x-1-\ln x.
\]
则
\[
g'(x)=1-\frac1x=\frac{x-1}{x}.
\]
$g$ 在 $(0,1)$ 上递减、在 $(1,+\infty)$ 上递增，所以
\[
g(x)\ge g(1)=0.
\]
故 $\ln x\le x-1$。

\basis 这里 $x=1$ 同时让 $\ln1=0$，也让导数的符号改变，是自然的锚点。

\reaction 构造差函数后，不是只求导。还要寻找一个函数值容易算、最好恰好为零的点，用单调性把整个区间接到这个点上。
\end{drill}

\begin{drill}{37. “恒成立”先把参数孤立}{项目导数题库恒成立题型改编}
若对任意 $x>0$，
\[
\ln x-ax\le0
\]
恒成立，求 $a$ 的取值范围。

\answer 原式等价于
\[
a\ge\frac{\ln x}{x}\quad(x>0).
\]
令 $g(x)=\dfrac{\ln x}{x}$，则
\[
g'(x)=\frac{1-\ln x}{x^2}.
\]
故 $g$ 在 $x=e$ 处取最大值 $\frac1e$。所以
\[
a\ge\frac1e.
\]

\basis “对任意 $x$，$a\ge g(x)$”要求 $a$ 压住 $g$ 的最大值。

\reaction 含参不等式恒成立，先把参数单独放一边。不要一上来对原式求导，先看它到底是在压最大值还是托最小值。
\end{drill}

\begin{drill}{38. “存在解”对应最值跨过零}{项目导数题库零点题型改编}
讨论方程
\[
x+\frac1x=a\qquad(x>0)
\]
何时有实数解。

\answer 令
\[
g(x)=x+\frac1x\quad(x>0).
\]
则
\[
g'(x)=1-\frac1{x^2}.
\]
$g$ 在 $(0,1)$ 上递减、在 $(1,+\infty)$ 上递增，最小值为 $g(1)=2$。因此方程有解当且仅当
\[
a\ge2.
\]

\basis 方程 $g(x)=a$ 有解，当且仅当 $a$ 落在 $g$ 的值域中。

\reaction 题目问“存在 $x$ 使某等式成立”时，把右边参数看成一条水平线；接下来问题是它能否碰到函数图像，也就是参数是否落入值域。
\end{drill}

\begin{drill}{39. 唯一零点要拆成存在和至多一个}{项目导数题库零点题型改编}
证明方程
\[
e^x+x-2=0
\]
有且只有一个实根。

\answer 令 $f(x)=e^x+x-2$。因为
\[
f'(x)=e^x+1>0,
\]
所以 $f$ 严格递增，至多有一个零点。又
\[
f(0)=-1<0,\qquad f(1)=e-1>0,
\]
由连续性知 $(0,1)$ 内至少有一个零点。故恰有一个。

\basis 单调性负责“至多一个”；端点异号加连续性负责“至少一个”。

\reaction 看见“有且只有一个零点”，强制拆成两半。别让零点存在定理替你偷偷证明唯一性，也别让单调性替你偷偷证明存在性。
\end{drill}

\begin{drill}{40. 导数符号能快速排除零点个数}{项目导数题库零点题型改编}
证明方程
\[
x^3+x+1=0
\]
只有一个实根。

\answer 令 $f(x)=x^3+x+1$。有
\[
f'(x)=3x^2+1>0,
\]
故 $f$ 严格递增，至多一个零点。又
\[
f(-1)=-1<0,\qquad f(0)=1>0,
\]
故至少一个零点。因此只有一个实根。

\basis 导数恒正说明图像只会上升一次，不可能与 $x$ 轴交两次。

\reaction 面对多项式零点个数，不必先设法分解。先问导数是否恒正或恒负；若是，唯一性部分已经完成。
\end{drill}

\begin{drill}{41. 导函数零点个数不等于原函数零点个数}{项目导数题库辨析题}
判断命题：“若 $f'(x)$ 没有零点，则 $f(x)$ 没有零点。”

\answer 错误。取
\[
f(x)=x.
\]
则 $f'(x)=1$ 没有零点，但 $f(0)=0$。

\basis $f'(x)=0$ 描述的是切线水平，不是函数值为零。一个是斜率问题，一个是纵坐标问题。

\reaction 题目把 $f(x)=0$ 和 $f'(x)=0$ 混在一起时，先问“它们分别等于什么”：前者是图像碰 $x$ 轴，后者是切线水平。
\end{drill}

\begin{drill}{42. 二阶导数研究的是一阶导数的增减}{项目导数题库改编}
已知
\[
f''(x)>0\quad(x\in\R),\qquad f'(1)=0.
\]
说明 $f'(x)$ 的正负，并判断 $x=1$ 处的极值类型。

\answer $f''(x)>0$ 说明 $f'$ 在 $\R$ 上严格递增。又 $f'(1)=0$，所以
\[
x<1\Rightarrow f'(x)<0,\qquad x>1\Rightarrow f'(x)>0.
\]
故 $f$ 在 $x=1$ 处取极小值。

\basis 二阶导数是导函数 $f'$ 的导数；$f''>0$ 不直接说 $f$ 递增，而是说 $f'$ 递增。

\reaction 看见 $f''$ 时先问“它在控制谁”。$f''$ 控制 $f'$ 的单调性，借助一个已知的 $f'$ 值才能进一步得到 $f$ 的增减。
\end{drill}

\begin{drill}{43. 导函数图像读的是原函数的增减}{人教B版导函数图像题型改编}
已知 $f'$ 在 $(0,1)$ 上为负，在 $(1,3)$ 上为正，在 $(3,+\infty)$ 上为负。写出 $f$ 的单调区间及极值点类型。

\answer $f$ 在 $(0,1)$ 上递减，在 $(1,3)$ 上递增，在 $(3,+\infty)$ 上递减。$x=1$ 处导数由负变正，是极小值点；$x=3$ 处由正变负，是极大值点。

\basis 导函数的符号就是原函数切线斜率的正负，从而决定原函数增减。

\reaction 读 $f'$ 的符号表时，把每一段直接翻译为“原函数升/降”；在分界点只看符号是怎样穿过去的。
\end{drill}

\begin{drill}{44. 导函数的极值不等于原函数的极值}{人教B版导函数图像题型改编}
已知 $f'$ 在 $x=2$ 处取得极小值 $-3$。能否据此断定 $f$ 在 $x=2$ 处取得极值？

\answer 不能。因为 $f'(2)=-3\ne0$，原函数在 $x=2$ 附近仍可能保持递减。$f'$ 的极值描述的是斜率变化得最快慢的位置，不是原函数的高低点。

\basis 原函数极值要看 $f'$ 在该点左右变号；仅知道 $f'$ 自己有极值，没有给出它是否过零。

\reaction 题目说谁的极值就盯谁：$f'$ 的极值是导函数的事；要判断 $f$ 的极值，仍然要看 $f'$ 的符号变化。
\end{drill}

\begin{drill}{45. 导数为负的区间要与定义域取交}{项目导数题库改编}
求函数
\[
f(x)=\frac{x}{x^2-1}
\]
的单调递减区间。

\answer 定义域为 $\R\setminus\{-1,1\}$，且
\[
f'(x)=\frac{(x^2-1)-2x^2}{(x^2-1)^2}
=-\frac{x^2+1}{(x^2-1)^2}<0
\]
在定义域内恒成立。因此递减区间为
\[
(-\infty,-1),\quad(-1,1),\quad(1,+\infty).
\]

\basis 导数在定义域内恒负，但定义域被两个无意义点切断；单调区间不能跨过函数没有定义的点。

\reaction 即使导数“恒正/恒负”，也先看定义域是否断开。单调区间是定义域中的连通区间，不能穿过空点。
\end{drill}

\begin{drill}{46. 极值点与极值是两个不同对象}{人教B版导数与极值题型改编}
函数
\[
f(x)=x^3-3x
\]
的极值点和极值分别是什么？

\answer
\[
f'(x)=3(x-1)(x+1).
\]
在 $x=-1$ 处导数由正变负，故极大值点为 $-1$，极大值为
\[
f(-1)=2.
\]
在 $x=1$ 处导数由负变正，故极小值点为 $1$，极小值为
\[
f(1)=-2.
\]

\basis 极值点是横坐标；极值是函数在该横坐标处取得的函数值。

\reaction 题目问“极值点”就回答 $x=$ 什么；问“极值”才把横坐标代回 $f(x)$。不要把 $x=1$ 和 $-2$ 混成同一种答案。
\end{drill}

\begin{drill}{47. 最值题的端点可能就是答案}{人教B版导数与最值题型改编}
求
\[
f(x)=\ln x-x
\]
在 $[1,4]$ 上的最大值。

\answer
\[
f'(x)=\frac1x-1.
\]
在 $(1,4)$ 内 $f'(x)<0$，故 $f$ 在 $[1,4]$ 上递减。因此最大值在左端点取得：
\[
f(1)=-1.
\]

\basis 即使没有内部极值点，闭区间最值仍存在，并可能直接在端点取得。

\reaction 求闭区间最值，不要把“求导数等于零”当成唯一动作。导数符号若已经说明单调性，就立刻回到端点比较。
\end{drill}

\begin{drill}{48. 参数出现在线性项时，导数常可转成二次函数最值}{项目导数题库改编}
若函数
\[
f(x)=x^3-3x^2+ax
\]
在 $\R$ 上单调递增，求 $a$ 的取值范围。

\answer
\[
f'(x)=3x^2-6x+a=3(x-1)^2+a-3.
\]
要对任意 $x$ 有 $f'(x)\ge0$，只需其最小值非负：
\[
a-3\ge0.
\]
故 $a\ge3$。

\basis 导函数是开口向上的二次函数；“对任意 $x$ 非负”等价于它的最小值不小于零。

\reaction 导函数是二次式且题目要求全体 $x$ 的单调性时，别急着判别式。先配方，看它的最小值是否被参数压到非负。
\end{drill}

\begin{drill}{49. “在某区间递增”不要求导函数在全体实数非负}{项目导数题库改编}
若函数
\[
f(x)=x^3-3x^2+ax
\]
在 $[2,+\infty)$ 上单调递增，求 $a$ 的取值范围。

\answer
\[
f'(x)=3(x-1)^2+a-3.
\]
在 $x\in[2,+\infty)$ 上，$(x-1)^2\ge1$，所以
\[
f'(x)\ge3+a-3=a.
\]
且 $x=2$ 时取到 $a$，故需且只需
\[
a\ge0.
\]

\basis 导函数的最小值要在题目指定区间上找，不是默认在全体实数上找。

\reaction 参数单调性题里，先圈出单调区间。导函数最值的取值范围由这个区间决定，不能无脑用全局顶点。
\end{drill}

\begin{drill}{50. 对数函数求导后的分母永远为正}{项目导数题库改编}
讨论函数
\[
f(x)=x-a\ln x\qquad(x>0)
\]
在 $(0,+\infty)$ 上的单调性，其中 $a>0$。

\answer
\[
f'(x)=1-\frac ax=\frac{x-a}{x}.
\]
因为 $x>0$，导数符号由 $x-a$ 决定。因此 $f$ 在
\[
(0,a)\text{ 上递减，在 }(a,+\infty)\text{ 上递增。}
\]

\basis 定义域 $x>0$ 使分母 $x$ 始终为正，所以分式导数的符号只需看分子。

\reaction 导函数出现分式时，不要立刻做复杂符号表。先用定义域判断分母恒正还是恒负，把符号问题缩小到分子。
\end{drill}

\begin{drill}{51. 导数比较可转化为函数增减比较}{人教B版导数概念题型改编}
已知 $f'(x)>g'(x)$ 对任意 $x\in(a,b)$ 成立。判断 $f(x)-g(x)$ 在 $(a,b)$ 上的单调性。

\answer 令
\[
h(x)=f(x)-g(x).
\]
则
\[
h'(x)=f'(x)-g'(x)>0.
\]
所以 $f(x)-g(x)$ 在 $(a,b)$ 上严格递增。

\basis 两个导数的大小关系可以移到同一个函数的导数正负问题上。

\reaction 比较两个函数或两个导数时，差函数是最自然的翻译：想比较 $f,g$，就考虑 $f-g$。
\end{drill}

\begin{drill}{52. 已知函数值和导数值时的切线近似}{人教B版导数应用题型改编}
已知 $f(1)=2,\ f'(1)=3$。用曲线在 $x=1$ 处的切线估计 $f(1.02)$。

\answer 切线为
\[
y-2=3(x-1).
\]
代入 $x=1.02$，得到
\[
f(1.02)\approx2+3\times0.02=2.06.
\]

\basis 在 $x_0$ 附近，曲线可用该点切线近似；切线方程就是
\[
f(x)\approx f(x_0)+f'(x_0)(x-x_0).
\]

\reaction 题目要求估计某个靠近已知点的函数值，先别硬代原函数。检查是否给了 $f(x_0)$、$f'(x_0)$；若给了，切线近似就是直接可用的信息。
\end{drill}

\end{document}
